\documentclass[11pt,a4paper]{article}
\usepackage[margin=23mm]{geometry}
\usepackage{amsmath,amssymb,amsthm}
\setlength{\emergencystretch}{2em}
\newtheorem{theorem}{Theorem}
\newcommand{\gen}[1]{\langle#1\rangle}
\newcommand{\Alt}{\operatorname{Alt}}
\title{The normalized orbit budget at odd degree}
\author{Human Compression project: L044 continuation}
\date{}
\begin{document}\maketitle
Retain a transitive connected string C-marking
$G=\gen{\rho_0,\ldots,\rho_{r-1}}\le\Alt(\Omega)$,
all ambient maximal omissions intransitive, $r\ge6$ and
$n=2r+1$. Assume the actual transitive six-point head
on $D=\{1,\ldots,6\}$ and the literal interface
\[
 \tau=\rho_3=(6\ b)\delta,\quad
 R=G_{\ge3}=\Alt(\Omega\setminus\{1,\ldots,5\}),\quad
 \sigma=\rho_4,\quad \sigma|_D=1.
\]
The odd involution $\delta$ fixes $D,b$. Write
$H=G_{\ge5}$ and $K=G_{>5}$, with $H$ fixing $D,b$.
Suppose every tail colour has two ambient-eligible
edges in one actual $H$-orbit. Normalize these
assignments by the complete whole-orbit transport.
In particular $K$ is intransitive on every orbit
with nonempty assigned set.

\begin{theorem}
Under these hypotheses $|\operatorname{Fix}(H)|\ge9$.
Thus $x=4$ supplies the local inequality
$|\operatorname{Fix}(G_{>x})|\ge2x+1$.
\end{theorem}
\begin{proof}
The acquired eight-fixed-point theorem excludes exactly
eight fixed points, and the standing interface gives
at least seven. It remains to exclude exactly seven.
Suppose this occurs. For each nontrivial $H$-orbit
$O_i$ put $d_i=|F_i|-|O_i|/2$, and let $p$ index
the orbit whose assigned set contains $5$.
The acquired normalized deficit theorem gives
\[
 \sum_i d_i=-2,\quad d_p\le0,\quad
 d_i\le-1\ (|O_i|\text{ even}),\quad
 d_i\le-\tfrac12\ (|O_i|\text{ odd})\quad(i\ne p).
                                                        \tag{1}
\]
Empty assignments have their exact deficit $-|O_i|/2$.
Every nontrivial internal action of $\delta$ is free,
because it centralizes the transitive group on that
orbit. On an orbit of size $2q$ the pair-forest bound
therefore gives deficit at most $-1$ for a nonempty
assignment. The same bound holds for an empty assignment;
if $q=1$ the assignment is necessarily empty, because
$\tau\in G_j$ already joins its two points for every
tail omission.

First $\delta$ preserves every nontrivial $H$-orbit.
Indeed, if it exchanges orbits of size $m$, transport
all assigned pairs of one orbit to the other by the
actual $\tau$. The complete selected-colour bound gives
combined deficit at most $-\lceil m/2\rceil$.
For $m=2$ or $3$ neither orbit contains two edges of
an involution, so both assignments are empty; their
combined deficit is $-m$. An exchange with odd $m$
therefore costs at least three deficit units and
contradicts (1). Every exchange of even orbits costs
at least two units and is itself even. Since $\delta$
fixes the seven fixed points and is odd, some further
odd internal action would be required; it costs at
least one more unit. The other deficits are nonpositive,
so even exchanges also contradict (1).

It follows that $\delta$ acts oddly on some preserved
orbit $O$ of size $2q$, with $q$ odd. The actual head
argument supplies
\[
       z=b^\sigma\ne b,\qquad \delta z=z,\qquad Kz=z,
       \quad z\notin D\cup\{b\}.                     \tag{2}
\]
In particular $z$ cannot belong to $O$.

We first dispose of every nonfirst saturated internal
orbit of size $2q\ge6$, irrespective of the parity of
$q$. Its deficit is $-1$ and its assigned set has
$q-1$ colours. The complete odd-degree row theorem
gives a terminal assigned label $t<r-1$ with
$|\operatorname{Fix}(G_{>t})|\ge2t+3$.
The nonempty remaining interval still has ambient
eligible pairs for each of its colours. Its complete
interval bound gives support at least $2(r-1-t)+1$.
Together these force $n\ge2r+2$, impossible.

Suppose now $q\ge3$ for the chosen odd orbit $O$.
An empty assignment would have deficit $-q\le-3$;
thus it is nonempty. If $O\ne O_p$, this excludes
deficit $-1$. Equation (1) leaves only deficit $-2$,
and all other deficits must vanish. Hence there
is precisely one further orbit $V=O_p$, of even
size $2s$ and deficit zero. Its two assigned
$5$-edges force $s\ge2$. This is exactly the
acquired nonfirst internal deficit-two theorem,
which excludes the configuration.

If $O=O_p$, its integral deficit is $-1$ or $-2$.
Deficit $-2$ would leave no further orbit by (1),
contradicting (2). With deficit $-1$, the remaining
orbits are either one even orbit $V$ of deficit
$-1$, or two odd orbits $V,W$ each of deficit
$-1/2$. The first alternative is exactly the
acquired first-internal theorem, including its
two-point residual, and is impossible.
In the second alternative both remaining assigned
sets are nonempty. The complete internal pair-path
equality gives an initial interval for $F_O$, and
the complementary terminal interval $I=F_V\cup F_W$
has all its eligible edges confined to $V\cup W$:
on $O$ its pair images are trivial, so its possible
edges there are already joined by $\tau$ in their
ambient omissions. Put $T=G_I$. Normalization makes
$K$ intransitive on each of $V,W$, and $T\le K$,
so $o_T(V\cup W)\ge4$. Complete localization gives
\[
 |V|+|W|-2=2|I|\le |V|+|W|-o_T(V\cup W)
                     \le |V|+|W|-4,
\]
a contradiction. Thus all possibilities with
$q\ge3$ have been excluded.

Finally let $q=1$, so $O=E$ is an empty assigned
pair of deficit $-1$. If $\delta$ fixes every other
orbit pointwise, it is literally one transposition,
and the acquired exterior-pair theorem applies with
all its natural-$R$, head and ambient-pair hypotheses.
It gives a contradiction. Otherwise let $W$ be a
further orbit on which $\delta$ acts nontrivially.
Its deficit is at most $-1$, hence exactly $-1$
by (1). If $|W|=2$, its assignment is also empty.
All remaining deficit is zero, so the only further
orbit is $O_p$, where $\delta$ must be trivial.
The two transpositions on $E,W$ would make $\delta$
even, a contradiction.

We may therefore suppose $|W|=2q'\ge4$. Its assigned
set is nonempty, because an empty one costs at least
two units. If $W=O_p$, (1) leaves no other orbit
besides $E,W$. Both are free under $\delta$, again
contradicting (2). If $W\ne O_p$, every remaining
deficit is zero and there is precisely one further
orbit $V=O_p$, of even size and deficit zero.
For $q'\ge3$, the saturated nonfirst exclusion above
applies. For $q'=2$, the literal budget is
\[
 |E|=2,\quad |W|=4,\quad
 F_E=\varnothing,\quad |F_W|=1,\quad
 F_V=\{5,\ldots,r-1\}\setminus F_W.
\]
The involution $\delta$ is free on $E,W$ and fixes
$V$ pointwise by (2) and centralization. This is
the full four-point budget excluded by the acquired
fixed-anchor corollary. It is impossible as well.
Every case with seven fixed points has now been
excluded. Together with the acquired eight-point
exclusion this proves the theorem.
\end{proof}

\paragraph{Attribution and boundary.}
The deficit, transport, internal pair-path, row-growth,
localization, fixed-anchor and residual orbit proofs
are acquired in full. This assembly supplies the
exhaustive seven-fixed-point case division and states
its local consequence. It does not produce an eligible
pair in one orbit when the two edges are split between
different orbits, and it does not assert that later
ambient generators preserve every locally fixed point.
Those assignment and attachment obligations remain
necessary for complete L044. No new rank induction,
computation, formal, original or question credit.
\end{document}
